%%%%%%%%%%%%%%%%%%%%%%%%%%%%%%%%%%%%%%%%%
% Beamer Presentation
% LaTeX Template
% Version 1.0 (10/11/12)
%
% This template has been downloaded from:
% http://www.LaTeXTemplates.com
%
% License:
% CC BY-NC-SA 3.0 (http://creativecommons.org/licenses/by-nc-sa/3.0/)
%
%%%%%%%%%%%%%%%%%%%%%%%%%%%%%%%%%%%%%%%%%

%----------------------------------------------------------------------------------------
%	PACKAGES AND THEMES
%----------------------------------------------------------------------------------------

\documentclass{beamer}
\usepackage{amsmath,amssymb,amsfonts}
\usepackage{graphicx}
\usepackage{tikz}
\usepackage{graphicx} % Allows including images
\usepackage{booktabs} % Allows the use of \toprule, \midrule and \bottomrule in tables
\usepackage{ytableau}

\mode<presentation> {

% The Beamer class comes with a number of default slide themes
% which change the colors and layouts of slides. Below this is a list
% of all the themes, uncomment each in turn to see what they look like.

%\usetheme{default}
%\usetheme{AnnArbor}
%\usetheme{Antibes}
%\usetheme{Bergen}
%\usetheme{Berkeley}
%\usetheme{Berlin}
%\usetheme{Boadilla}
%\usetheme{CambridgeUS}
%\usetheme{Copenhagen}
%\usetheme{Darmstadt}
%\usetheme{Dresden}
%\usetheme{Frankfurt}
%\usetheme{Goettingen}
%\usetheme{Hannover}
%\usetheme{Ilmenau}
%\usetheme{JuanLesPins}
%\usetheme{Luebeck}
\usetheme{Madrid}
%\usetheme{Malmoe}
%\usetheme{Marburg}
%\usetheme{Montpellier}
%\usetheme{PaloAlto}
%\usetheme{Pittsburgh}
%\usetheme{Rochester}
%\usetheme{Singapore}
%\usetheme{Szeged}
%\usetheme{Warsaw}

% As well as themes, the Beamer class has a number of color themes
% for any slide theme. Uncomment each of these in turn to see how it
% changes the colors of your current slide theme.

%\usecolortheme{albatross}
%\usecolortheme{beaver}
%\usecolortheme{beetle}
%\usecolortheme{crane}
%\usecolortheme{dolphin}
%\usecolortheme{dove}
%\usecolortheme{fly}
%\usecolortheme{lily}
%\usecolortheme{orchid}
%\usecolortheme{rose}
%\usecolortheme{seagull}
%\usecolortheme{seahorse}
%\usecolortheme{whale}
%\usecolortheme{wolverine}

%\setbeamertemplate{footline} % To remove the footer line in all slides uncomment this line
%\setbeamertemplate{footline}[page number] % To replace the footer line in all slides with a simple slide count uncomment this line

%\setbeamertemplate{navigation symbols}{} % To remove the navigation symbols from the bottom of all slides uncomment this line
}

\newcommand{\youngone}[1]{\begin{tikzpicture}[scale=0.4]
\foreach \i in {1,...,#1} \draw (\i-1,0) rectangle (\i,1);
\end{tikzpicture}}
\newcommand{\youngtwo}[2]{\begin{tikzpicture}[scale=0.4]
\foreach \i in {1,...,#1} \draw (\i-1,0) rectangle (\i,1);
\foreach \i in {1,...,#2} \draw (\i-1,-1) rectangle (\i,0);
\end{tikzpicture}}
\newcommand{\youngthree}[3]{\begin{tikzpicture}[scale=0.4]
\foreach \i in {1,...,#1} \draw (\i-1,0) rectangle (\i,1);
\foreach \i in {1,...,#2} \draw (\i-1,-1) rectangle (\i,0);
\foreach \i in {1,...,#3} \draw (\i-1,-2) rectangle (\i,-1);
\end{tikzpicture}}

\setlength{\unitlength}{1mm}

% 定义一个带箭头的“从 A 指向 B”的小宏
% 语法：\link(x1,y1)(x2,y2)
\newcommand{\link}[4]{\begin{picture}(0,0)
  \put(#1,#2){\vector(#3,#4){10}} % 10 是长度，可自行调
\end{picture}}
\renewcommand{\arraystretch}{1.5}
%----------------------------------------------------------------------------------------
%	TITLE PAGE
%----------------------------------------------------------------------------------------

\title[Short title]{Convergence of the DFP Algorithm} % The short title appears at the bottom of every slide, the full title is only on the title page

\author{Jinwen Huang} % Your name
\institute[sysu] % Your institution as it will appear on the bottom of every slide, may be shorthand to save space
{
Sun Yat-sen University \\ % Your institution for the title page
\medskip
\textit{huangjw267@mail2.sysu.edu.cn} % Your email address
}
\date{\today} % Date, can be changed to a custom date

\begin{document}

\begin{frame}
\titlepage % Print the title page as the first slide
\end{frame}

\begin{frame}
\frametitle{Overview} % Table of contents slide, comment this block out to remove it
\tableofcontents % Throughout your presentation, if you choose to use \section{} and \subsection{} commands, these will automatically be printed on this slide as an overview of your presentation
\end{frame}

%----------------------------------------------------------------------------------------
%	PRESENTATION SLIDES
%----------------------------------------------------------------------------------------

%------------------------------------------------
\section{The DFP Algorithm} 
%------------------------------------------------
\subsection{The DFP Algorithm}
\begin{frame}
\frametitle{The DFP Algorithm}
    The DFP Algorithm is the first quasi-Newton method for unconstrained optimization \[\min _{x \in \mathbb{R}^{n}} f(x) . \tag{1.1}\]
    The algorithm is as follows.
\end{frame}

\begin{frame}
\frametitle{The DFP Algorithm}
Step 0. Given $x_1 \in \mathbb{R}^n ; B_1 \in \mathbb{R}^{n \times n}$ positive definite;$k:=1.$\\
Step 1. Compute $g_k=\nabla f\left(x_k\right)$;\\
if $g_k=0$ then stop;\\
set $d_k=-B_k^{-1} g_k$;\\
Step 2. Carry out a line search along $d_k$, getting $\alpha_k>0$;\\
set $x_{k+1}=x_k+\alpha_k d_k$;\\
Step 3. Set\begin{equation*}
    B_{k+1}=B_k-\frac{B_k s_k y_k^T+y_k s_k^T B_k}{s_k^T y_k}+\left(1+\frac{s_k^T B_k s_k}{s_k^T y_k}\right) \frac{y_k y_k^T}{s_k^T y_k}, \tag{1.2}
\end{equation*}
where\begin{align*}
    s_k=\alpha_k d_k \tag{1.3}\\
    y_k=g_{k+1}-g_k  \tag{1.4}
\end{align*}
Step 4.k:=k+1 \text {; go to Step } 1.

\end{frame}

\begin{frame}
    \frametitle{The DFP Algorithm}
    If we replace(1.2)with\begin{equation*}
B_{k+1}=B_{k}-\frac{B_{k} s_{k} s_{k}^{T} B_{k}}{s_{k}^{T} B_{k} s_{k}}+\frac{y_{k} y_{k}^{T}}{s_{k}^{T} y_{k}}
\end{equation*}
    then we get the BFGS algorithm.
    
\end{frame}



%------------------------------------------------
\section{The Wolfe Line Search}
%------------------------------------------------
\subsection{The Wolfe Line Search}

\begin{frame}{The Wolfe Line Search}
The line search in Step 2 requires step lengths $d_k$ to satisfy certain line search conditions.\\
If exact line search is used, $\alpha_k$ satisfies\begin{equation*}
    f\left(x_k+\alpha_k d_k\right)=\min f\left(x_k+\alpha d_k\right) .\tag{2.1}
\end{equation*}

For inexact line search, the Wolfe conditions usually serves as an alternative for.\\

The Wolfe Line Search consists of two conditions, the Armijo Condition that ensures sufficient decrease and the Curvature Condition ensures the steps size to be "good" enough.
\end{frame}

\begin{frame}
    \frametitle{Armijo Condition}
The Armijo condition is a sufficient decrease condition that ensures that the step size chosen in the line search leads to a sufficient decrease in the objective function.
The formula for the Armijo condition is given by:
\[f(x_k + \alpha d_k) \leq f(x_k) + c_1 \alpha \nabla f(x_k)^T d_k,\tag{2.2} \]
where $f$ is the objective function, $x_k$ is the current point, $d_k$ is the search direction, $\alpha$ is the step size, and $c_1$ is a constant in the range $(0, 1)$.  The Armijo condition ensures that the step size chosen in the line search leads to a sufficient decrease in the objective function, which is important for the convergence of optimization algorithms.   

\end{frame}

\begin{frame}
    \frametitle{Curvature Condition}
The curvature condition is a condition that ensures that the step size chosen in the line search is not too small. 
It is given by the formula:
\[\nabla f(x_k + \alpha d_k)^T d_k \geq c_2 \nabla f(x_k)^T d_k,\tag{2.3} \]
where $c_2$ is a constant in the range $(c_1, 1)$. The curvature condition ensures that the step size chosen in the line search is not too small, which can lead to slow convergence or even divergence of optimization algorithms.


\end{frame}



%------------------------------------------------



\section{Advancements}

\begin{frame}
    \frametitle{Advancements}
    \begin{itemize}
    \item Powell\cite{Po1} showed that if $f(x)$ is a convex function and if line searches are exact in all iterations, then the DFP algorithm will either stop at a global minimum or generate a sequence of iterates that converges to a global minimum. 
    \item For inexact line search, Powell\cite{Po2} showed that the BFGS method is globally convergent. 
    \item Byrd, Nocedal and Yuan\cite{BNY} showed the global convergence of all the methods in the convex Broden family except the DFP method.

\end{itemize}
\end{frame}

\begin{frame}
    \frametitle{Advancements}
\begin{itemize}
    \item In the 1990s, Yuan\cite{Yuan}(improved by Yin and Han\cite{YH} three years later), Xu and Liu\cite{XL} gave sufficient conditions for the global convergence of the DFP method with objective function being uniformly convex respectively.
    \item In 2002, Liu, Han and Xu\cite{LHX} proved that convergence of DFP algorithm with Wolfe line search holds for quadratic uniformly convex functions.
    \item In 2024, Luo, Wu and He\cite{LWH} derived a new DFP correction formula and proposed a new class of DFP algorithm under the strong Wolfe condition, improving the problem of linesearch failure due to accuracy and proving global convergence of the improved algorithm under certain assumptions.
\end{itemize}
\end{frame}



\section{Some sufficient conditions for the global convergence of the DFP algortihm}

\begin{frame}
    \frametitle{Some sufficient conditions for the global convergence of the DFP algortihm}
    In this section, we will talk about some sufficient conditions for the global convergence of the DFP Algorithm given by Yuan\cite{Yuan}, Xu and Liu\cite{XL}, Yin and Han\cite{YH}.\\
    All these works were established under the following hypothesis (H):
    \begin{enumerate}
        \item f(x) is twice continuously differentiable and uniformly convex. Namely, there exists a constant $c_{3} \geq 0$ such that
    \begin{equation*}
        d^{T} \nabla^{2} f(x) d \geqslant c_{3}\|d\|_{2}^{2}, \quad \forall d, x \in \mathbb{R}^{n} . \tag{3.1}
    \end{equation*}\\
    \item The Wolfe Conditions are satisfied in all iterations.
    \end{enumerate}
\end{frame}


\subsection{Yuan's work in the 1990s}

\begin{frame}
    \frametitle{Yuan's work in the 1990}
        Using the well-known Trace Relation\begin{equation}
            \operatorname{tr}\left(B_{k+1}\right)=\operatorname{tr}\left(B_{k}\right)-2 \frac{s_{k}^{T} B_{k} y_{k}}{s_{k}^{T} y_{k}}+\left(1+\frac{s_{k}^{T} B_{k} s_{k}}{s_{k}^{T} y_{k}}\right) \frac{\left\|y_{k}\right\|_{2}^{2}}{s_{k}^{T} y_{k}},  \tag{3.2}
        \end{equation}
        Yuan obtained the following result:
        \begin{theorem}[Yuan,2.1]
            If $\left\{\alpha_{k}\right\}$ are bounded away from zero, and if $\operatorname{Tr}\left(B_{k}\right)$ are monotonically increasing, the DFP method converges to the unique solution.
        \end{theorem}
    

\end{frame}

\begin{frame}
    \frametitle{Yuan's work in the 1990}
    Subsequently, Yuan got the following:
    \begin{theorem}[Yuan,2.2]
       Assume that the hypothesis (H) holds. If the sequence $x_{k}$ generated by the DFP algorithm satisfies
\begin{equation*}
\left\|g_{k+1}\right\|_{2} \leqslant\left\|g_{k}\right\|_{2} \tag{2.21}
\end{equation*}
for all large $k$, then $x_{k}$ converges to the unique minimum of $f(x)$.
    \end{theorem}
    
\end{frame}

\begin{frame}
    \frametitle{Yuan's work in the 1990}
    In addition, we have
        \begin{theorem}[Yuan,2.3]
    If
\begin{equation*}
\sum_{k=1}^{x}| | s_{k}| |<\infty . \tag{2.43}
\end{equation*}
then the DFP algorithm converges to the solution.
    \end{theorem}
    \begin{theorem}[Yuan,2.4]
        Assume that the conditions in Hypothesis (H) hold. Let the sequence $\{x_{k}, k=1,2, \cdots\}$ be generated by the DFP method. If
\begin{equation*}
y_{k}^{T} B_{k} g_{k} \leqslant 0 \tag{2.47}
\end{equation*}
for all $k$, then $\left\{x_{k}\right\}$ either terminates at the unique minimum $x^{*}$ of $f(x)$ or converges to $x^{*}$.
    \end{theorem}
\end{frame}



\subsection{Xu and Liu's work}

\begin{frame}
\frametitle{Xu and Liu's work}
In this essay, Xu and Liu established the following sufficient conditions for the global convergence of the DFP Algorithm.
\begin{theorem}[Xu and Liu,2.3]
    Assume the conditions in Hypothesis (H) hold. If \begin{equation*}
        y_k^{\mathrm{T}} B_k^2 g_k \leq 0 \tag{2.10}
    \end{equation*}
holds for all large $k$, then the sequence $\left\{x_k, k=1,2, \cdots\right\}$ generated by the DFP method either terminates at the unique minimum $x^*$ of $f(x)$ or converges to $x^*$.
\end{theorem}
\end{frame}

\begin{frame}
\frametitle{Xu and Liu's work}
\begin{theorem}[Xu and Liu,2.4]
The condition (2.10) in Theorem (Xu and Liu,2.3)can be replaced by \begin{equation*}
    \sum_{i \in I_k}\left(y_i^{\mathrm{T}} B_i^2 g_i\right)^2 \leq C_2 \sum_{i=1}^k \frac{1}{\left\|d_i\right\|_2^3},
\end{equation*}
where $I_k=\left\{i: y_i^{\mathbf{T}} B_i^2 g_i>0, i<k\right\}, C_2>0$.
\end{theorem}

\end{frame}



\subsection{Yin and Han's Improvement}

\begin{frame}
\frametitle{Yin and Han's Improvement}
In this essay, Yin and Han established the following theorem:
\begin{theorem}[Yin and Han,2.3]
Assume that the conditions in Hypothesis (H) hold. Let the sequence $\{x_{k}, k=1,2, \cdots\}$ be generated by the DFP method. If
\begin{equation*}
    \operatorname{tr}\left(B_k\right) \leq c_6 \sqrt{\sum_{i=1}^k \frac{1}{\left\|d_i\right\|^2}}, \quad \forall k \geq 1, \tag{2.14}
\end{equation*}
where $c_6>0$ constant,then $\left\{x_k\right\}$ terminates at the unique minimum $x^{*}$ of $f(x)$ or converges to $x^{*}$.
\end{theorem}
    
\end{frame}

\begin{frame}
\frametitle{Yin and Han's Improvement}
Then we get the following theorem:
\begin{theorem}[Yin and Han,2.4]
    Under the hypothesis in theorem (Yin and Han,2.3), if there exists a constant $c_7>0$ such that \begin{equation*}
        \operatorname{Tr}\left(B_k\right) \leq c_7 k, \quad \forall k \geq 1,
    \end{equation*}
then $\left\{x_k\right\}$ terminates at the unique minimum $x^{*}$ of $f(x)$ or converges to $x^{*}$.
\end{theorem}
    From the theorems above, we can get some sufficient conditions for the global convergence of the DFP Algorithm.
\end{frame}

\begin{frame}
\frametitle{Yin and Han's Improvement}
\begin{theorem}[Yin and Han,3.1]
Assume that the conditions in Hypothesis (H) hold. Let the sequence $\{x_{k}, k=1,2, \cdots\}$ be generated by the DFP method. If\begin{equation*}
    \sum_{i \in I_k} \frac{g_i^T y_i}{d_i^T y_i} \leq c \sqrt{\sum_{i=1}^k \frac{1}{\left\|d_i\right\|^2}}, \quad \forall k \geq 1,
\end{equation*}
where $c>0$ constant,$I_k=\left\{i \mid 1 \leq i \leq k, g_i^T y_i>0\right\}$,then $\left\{x_k\right\}$ terminates at the unique minimum $x^{*}$ of $f(x)$ or converges to $x^{*}$.
    
\end{theorem}
    
\end{frame}

\begin{frame}
\frametitle{Yin and Han's Improvement}
    \begin{theorem}[Yin and Han,3.2]
Assume that the conditions in Hypothesis (H) hold. Let the sequence $\{x_{k}, k=1,2, \cdots\}$ be generated by the DFP method. If \begin{equation*}
    g_k^T y_k \leq 0, \quad \forall k \geq 1,
\end{equation*}
then $\left\{x_k\right\}$ terminates at the unique minimum $x^{*}$ of $f(x)$ or converges to $x^{*}$.
\end{theorem}

Note that 
\begin{align*}
g_k^T y_k & =g_k^T\left(g_{k+1}-g_k\right)=g_k^T g_{k+1}-\left\|g_k\right\|^2 \leq \frac{1}{2}\left(\left\|g_k\right\|^2+\|g_{k+1} \|^2\right)-\left\|g_k\right\|^2 \\
& =\frac{1}{2}\left(\left\|g_{k+1}\right\|^2-\left\|g_k\right\|^2\right) \leq 0, \quad \forall k \geq 1
\end{align*}
then from theorem (Yin and Han,3.2),we obtain theorem (Yuan,2.2), indicating that the conditions in theorem (Yuan,2.2) are stronger than those in theorem(Yin and Han,3.1).
\end{frame}

\begin{frame}
\frametitle{Yin and Han's Improvement} 
\begin{theorem}[Yin and Han,3.4]
Assume that the conditions in Hypothesis (H) hold. Let the sequence $\{x_{k}, k=1,2, \cdots\}$ be generated by the DFP method. If \begin{equation*}
    \sum_{k \in J} g_k^T y_k<+\infty, \tag{3.10}
\end{equation*}
where $J=\left\{k \mid g_k^T y_k>0, k=1,2, \cdots\right\}$,then $\left\{x_k\right\}$ terminates at the unique minimum $x^{*}$ of $f(x)$.
\end{theorem}
If we replace (3.10) with
\begin{align*}
    \sum_{k=1}^{\infty}\left|g_k^T y_k\right|<+\infty, \tag{3.12}
\end{align*}
theorem (Yin and Han,3.4) also holds.
\end{frame}

\begin{frame}
\frametitle{Yin and Han's Improvement} 
Note that
\begin{align*}
    \left|g_k^T y_k\right| \leq\left\|g_k\right\|\left\|y_k\right\| \leq \eta\left\|y_k\right\| \leq \eta c_4\left\|s_k\right\|,
\end{align*}
then from theorem (Yin and Han,3.4),we obtain theorem (Yuan,2.3), indicating that the conditions in theorem (Yuan,2.3) are stronger than  (3.12).\\
Additionally, we can also obtain (Yuan,2.1) by (Yin and Han,3.4).
\end{frame}

\begin{frame}
\frametitle{Yin and Han's Improvement} 
\begin{theorem}
    Assume that the conditions in Hypothesis (H) hold. If there exists a constant $b>0$ such that \begin{equation*}
    \left\|g_k\right\| \leq b\left\|d_k\right\|, \quad \forall k \geq 1,
\end{equation*}
then $\left\{x_k\right\}$ terminates after finite steps or converges to the unique minimum $x^{*}$ of $f(x)$.
\end{theorem}
\end{frame}

\begin{frame}
\frametitle{Yin and Han's Improvement} 
Under the hypothesis (H), Yuan\cite{Yuan} gave an sufficient condition
\begin{align*}
    y_k^T B_k g_k \leq 0, \quad \forall k \geq 1.
\end{align*}
Xu and Liu\cite{XL} suggested that
\begin{align*}
    y_k^T B_k^2 g_k \leq 0, \quad \forall k \geq 1.
\end{align*}
Yin and Han\cite{YH} showed that
\begin{align*}
    y_k^T B_k^0 g_k \leq 0, \quad \forall k \geq 1.
\end{align*}
One may propose that if there exists an integer $m\geq0$ such that
\begin{align*}
    y_k^T B_k^m g_k \leq 0, \quad \forall k \geq 1
\end{align*}
then $\left\{x_k\right\}$ terminates after finite steps or converges to the unique minimum $x^{*}$ of $f(x)$. However, whether this proposition holds is still unknown.
\end{frame}



\section{Summary}
\subsection{}

\begin{frame}
In this talk, we discover the origin and recent progress in the global convergence of the DFP algorithm.\\
Then we review some essays giving sufficient conditions of convergence under some hypotheis.((i)the functions are twice continuously differentiable and uniformly convex (ii) Wolfe conditions are satisfied in all iterations)\\
Finally, we make some comparisons of the results, and point out an open problem to be solved.

\end{frame}
%------------------------------------------------

\begin{frame}
\frametitle{References}
\footnotesize{
\begin{thebibliography}{99} % Beamer does not support BibTeX so references must be inserted manually as below
\bibitem[1]{Da} W. C. Davidon. 
\newblock Variable metric method for minimization.
\newblock \textit{Argonne National Laboratories Report ANL-5990}, 1959.
\end{thebibliography}

\begin{thebibliography}{99} % Beamer does not support BibTeX so references must be inserted manually as below
\bibitem[2]{FP} R. Fletcher and M. J. D. Powell.
\newblock    A rapidly convergent descent method for minimization.
\newblock   \textit{The Computer Journal}, 6 (1963), 163–168.
\end{thebibliography}

\begin{thebibliography}{99} % Beamer does not support BibTeX so references must be inserted manually as below
\bibitem[3]{Po1} M. J. D. Powell.
\newblock    On the convergence of variable metric algorithm.
\newblock    \textit{J. Inst. Math. Appl.}, (1971), 7:21.
\end{thebibliography}

\begin{thebibliography}{99} % Beamer does not support BibTeX so references must be inserted manually as below
\bibitem[4]{Po2} M. J. D. Powell.
\newblock    Some Global convergence properties of a variable metric algorithm for minimization without exact line searches.
\newblock    \textit{Nonlinear Programming, SIAM-AMS Proceedings vol. IX}, 1976, 53–72.
\end{thebibliography}

\begin{thebibliography}{99} % Beamer does not support BibTeX so references must be inserted manually as below
\bibitem[5]{BNY} R. H. Byrd, J. Nocedal and Y. Yuan.
\newblock    Global convergence of a class of variable metric algorithms.
\newblock   \textit{SIAM J. Numer. ANal.}, 1987, 24:1171.
\end{thebibliography}
}

\end{frame}

\begin{frame}

\frametitle{References}
\footnotesize{

\begin{thebibliography}{99}
\bibitem[6]{Yuan} Y. Yuan.
\newblock    On the convergence of the DFP algorithm.
\newblock    Technical Report, Computing Center, Academica Sinica, Beijing, China, 1994.
\end{thebibliography}

\begin{thebibliography}{99}
\bibitem[7]{YH} H. Yin and J. Han.
\newblock    Some sufficient conditions for the global convergence of the DFP algorithm.
\newblock   \textit{Acta Math. Appl. Sin. 21}, No.2, 179-186, 1998.
\end{thebibliography}

\begin{thebibliography}{99}
\bibitem[8]{XL} D. Xu and G. Liu.
\newblock    A new sufficient condition for the convergence of DFP algorithm with Wolfe line search.
\newblock    \textit{Systems Science and Mathematical Sciences}, 9(1996), 259-269.    
\end{thebibliography}

\begin{thebibliography}{99}
\bibitem[9]{LHX} G. Liu, J. Han and D. Xu.
\newblock    A note on the convergence of the DFP algoithm on quadratic uniformly convex functions.
\newblock    \textit{Optimization 51}, No.2, 339-352, 2002.    
\end{thebibliography}

\begin{thebibliography}{99}
\bibitem[10]{LWH} W. Luo, Z. Wu and S. He.
\newblock    An Improved Quasi-Newtonian Algorithm
\newblock    \textit{Journal of Chengdu University of Information Technology}, 2024, 39(03): 374-381. 
\end{thebibliography}
}

\end{frame}

%------------------------------------------------

\begin{frame}
\Huge{\centerline{Thanks For Listening}}
\end{frame}

%----------------------------------------------------------------------------------------

\end{document}